%% ch11_human_side.tex -- Intelligence Engineering, chapter 11 "The Human Side
%% of Machine Learning".
%% Spine transferred from Huyen, Designing Machine Learning Systems, chapter
%% 11. The subchapters and the frame titles under them are the book's own
%% section and subsection headings, in the book's order.
%%
%% STATE: titles and TEMPLATE layout only. No body text. Every frame carries
%% the layout it is meant to be filled into, and a % TEMPLATE comment saying
%% why that layout and what belongs in each slot. Filling them is the next pass.
%%
%% Fragment of intelligence_engineering.tex. One chapter is one lecture and
%% gets exactly ONE \lecturepage, at the top. Inside it every \subsection
%% opens on the subchapter divider, which the master emits automatically via
%% \AtBeginSubsection; do not write those divider frames by hand.

\begin{refsection}

\lecturepage{XI}{The Human Side of Machine Learning}{}
\section[Human Side]{The Human Side of Machine Learning}

% ----------------------------------------------------------------------
\subsection{User Experience}

% TEMPLATE: withmargin, the chapter's workhorse and the right opener for a
% requirement: one claim plus the aside that makes it concrete. Main column
% takes what consistency means to a user and the tension the book sets it
% against. Margin holds the book's worked case, the travel site whose model
% overrode the filters a user had set, and the verdict its designers reached.
\begin{frame}{Ensuring User Experience Consistency}
  \begin{withmargin}
    \begin{tdmain}
    \end{tdmain}
    \begin{tdmargin}
    \end{tdmargin}
  \end{withmargin}
  \slidesources{Huyen2022}
\end{frame}

% TEMPLATE: side image hero, the first of this chapter's two. The book makes
% this heading concrete with a figure of several predictions offered at once,
% so the strip takes that figure and the sidebody takes why a mostly correct
% prediction is worthless to a user who cannot check it, and what showing many
% predictions buys instead. The gray placeholder stays until the asset exists.
\begin{frame}{Combatting ``Mostly Correct'' Predictions}
  \phimgside[0.33]{combatting_mostly_correct_predictions.png}{Mostly correct predictions}
  \begin{sidebody}
  \end{sidebody}
  \slidesources{Huyen2022}
\end{frame}

% TEMPLATE: tddef plus a takeaway, and deliberately not a third layout of the
% same kind in a row. The heading is the book's own term for a backup route,
% so the term box is the slide: it takes the definition and the rule that
% decides when the fallback fires. The takeaway keeps the speed against
% accuracy tradeoff the book closes the heading with.
\begin{frame}{Smooth Failing}
  \begin{tddef}{Smooth failing}
  \end{tddef}
  \begin{takeaway}
  \end{takeaway}
  \slidesources{Huyen2022}
\end{frame}

% ----------------------------------------------------------------------
\subsection{Team Structure}

% TEMPLATE: withmargin. The heading is about who has to be in the room, so the
% main column takes the subject matter experts the book names and the stages it
% says they must reach beyond labeling. Margin holds the friction the book
% reports once those experts are asked to work inside engineers' tooling.
\begin{frame}{Cross-functional Teams Collaboration}
  \begin{withmargin}
    \begin{tdmain}
    \end{tdmain}
    \begin{tdmargin}
    \end{tdmargin}
  \end{withmargin}
  \slidesources{Huyen2022}
\end{frame}

% TEMPLATE: two block contrast, levelled by the equal height group. This frame
% opens the two deep headings under it, so it stays a comparison and not a
% detail layout. The headers are the two approaches the children name; each box
% takes that approach in a line, and the children carry the workflow itself.
\begin{frame}{End-to-End Data Scientists}
  \begin{columns}[T,onlytextwidth]
    \begin{column}{0.47\textwidth}
      \begin{tdblockt}[equal height group=e2e]{Approach 1}
      \end{tdblockt}
    \end{column}
    \begin{column}{0.47\textwidth}
      \begin{tdblockt}[equal height group=e2e]{Approach 2}
      \end{tdblockt}
    \end{column}
  \end{columns}
  \slidesources{Huyen2022}
\end{frame}

% TEMPLATE: side image hero, the second and last of this chapter's two. The
% book draws this approach as a handoff, the data scientists on one side and
% the team that owns production on the other, so the strip takes that diagram
% and the sidebody takes the communication overhead the book charges it with.
\begin{frame}{Approach 1: Have a separate team to manage production}
  \phimgside[0.33]{approach_1_have_a_separate_team_to_manage_production.png}{A separate team to manage production}
  \begin{sidebody}
  \end{sidebody}
  \slidesources{Huyen2022}
\end{frame}

% TEMPLATE: withmargin. The heading is one claim about ownership, so the main
% column takes what owning the entire process means step by step and the
% infrastructure it presupposes. Margin holds the named practice the book cites
% as its case and the objection it quotes from the data scientists themselves.
\begin{frame}{Approach 2: Data scientists own the entire process}
  \begin{withmargin}
    \begin{tdmain}
    \end{tdmain}
    \begin{tdmargin}
    \end{tdmargin}
  \end{withmargin}
  \slidesources{Huyen2022}
\end{frame}

% ----------------------------------------------------------------------
\subsection{Responsible AI}

% TEMPLATE: booktabs comparison table, and the opener for the two case studies
% under it. One column per case study named in the children, so the slide reads
% as a pair rather than as two sentences. Row labels are left empty on purpose:
% they are the dimensions the next pass fills, what the system did and who
% carried the cost.
\begin{frame}{Irresponsible AI: Case Studies}
  \footnotesize
  \renewcommand{\arraystretch}{1.35}
  \begin{tabular}{@{}%
    >{\raggedright\arraybackslash}p{0.22\textwidth}%
    >{\raggedright\arraybackslash}p{0.35\textwidth}%
    >{\raggedright\arraybackslash}p{0.35\textwidth}@{}}
    \toprule
     & \textbf{Case study I} & \textbf{Case study II}\\
    \midrule
     & & \\
     & & \\
     & & \\
     & & \\
    \bottomrule
  \end{tabular}
  \slidesources{Huyen2022}
\end{frame}

% TEMPLATE: withmargin. The case is an argument from numbers, so the main
% column runs what the grader was asked to do and what it graded on instead.
% Margin takes the book's figures on where the downgrades landed, which is the
% evidence the whole case rests on.
\begin{frame}{Case study I: Automated grader's biases}
  \begin{withmargin}
    \begin{tdmain}
    \end{tdmain}
    \begin{tdmargin}
    \end{tdmargin}
  \end{withmargin}
  \slidesources{Huyen2022}
\end{frame}

% TEMPLATE: one plain box plus a takeaway, and not a second withmargin in a
% row. The heading already carries the claim, so the box needs no header over
% it: it takes the released dataset, the fields left inside it and the
% reidentification that followed. The takeaway is the book's verdict on what
% the word anonymized actually buys.
\begin{frame}{Case study II: The danger of ``anonymized'' data}
  \begin{tdblock}
  \end{tdblock}
  \begin{takeaway}
  \end{takeaway}
  \slidesources{Huyen2022}
\end{frame}

% TEMPLATE: withmargin, opener for the seven deep headings under it. Seven
% peers do not grid, so this frame stays one claim about why a framework is
% needed at all and the children carry one recommendation each. Margin holds
% the book's own caveat, that this is a starting point and not a standard.
\begin{frame}{A Framework for Responsible AI}
  \begin{withmargin}
    \begin{tdmain}
    \end{tdmain}
    \begin{tdmargin}
    \end{tdmargin}
  \end{withmargin}
  \slidesources{Huyen2022}
\end{frame}

% TEMPLATE: five block grid, three across then two, levelled by two equal
% height groups. The book answers this heading with peer sources of bias spread
% along the pipeline, none subordinate to another, so one box each rather than
% a bullet list. Boxes stay untitled: the heading names no source, the names
% arrive with the text. Trim or add a box to the count the book gives.
\begin{frame}{Discover sources for model biases}
  \begin{columns}[T,onlytextwidth]
    \begin{column}{0.31\textwidth}
      \begin{tdblock}[equal height group=bias1]
      \end{tdblock}
    \end{column}
    \begin{column}{0.31\textwidth}
      \begin{tdblock}[equal height group=bias1]
      \end{tdblock}
    \end{column}
    \begin{column}{0.31\textwidth}
      \begin{tdblock}[equal height group=bias1]
      \end{tdblock}
    \end{column}
  \end{columns}
  \vskip0.4em
  \begin{columns}[T,onlytextwidth]
    \begin{column}{0.47\textwidth}
      \begin{tdblock}[equal height group=bias2]
      \end{tdblock}
    \end{column}
    \begin{column}{0.47\textwidth}
      \begin{tdblock}[equal height group=bias2]
      \end{tdblock}
    \end{column}
  \end{columns}
  \slidesources{Huyen2022}
\end{frame}

% TEMPLATE: withmargin. The heading names a limitation and not a set of peers,
% so the main column takes what the data on its own cannot settle and what the
% book asks for in its place. Margin holds the case it argues from and the
% people it says the data never reached.
\begin{frame}{Understand the limitations of the data-driven approach}
  \begin{withmargin}
    \begin{tdmain}
    \end{tdmain}
    \begin{tdmargin}
    \end{tdmargin}
  \end{withmargin}
  \slidesources{Huyen2022}
\end{frame}

% TEMPLATE: booktabs table, both column headers taken from the heading's own
% words. A tradeoff is a pair, so one row per desideratum against what it
% costs reads as a table and not as prose. Rows stay empty for the pairs the
% book names and the argument it makes for each.
\begin{frame}{Understand the trade-offs between different desiderata}
  \footnotesize
  \renewcommand{\arraystretch}{1.35}
  \begin{tabular}{@{}%
    >{\raggedright\arraybackslash}p{0.45\textwidth}%
    >{\raggedright\arraybackslash}p{0.45\textwidth}@{}}
    \toprule
    \textbf{Desideratum} & \textbf{Trade-off}\\
    \midrule
     & \\
     & \\
     & \\
    \bottomrule
  \end{tabular}
  \slidesources{Huyen2022}
\end{frame}

% TEMPLATE: one titled box plus a takeaway. The heading is a short imperative,
% so it becomes the box header itself and the box takes the reason behind it,
% what a bias costs once everything else is already built on top of it. The
% takeaway carries the book's line on the cheapest moment to act.
\begin{frame}{Act early}
  \begin{tdblockt}{Act early}
  \end{tdblockt}
  \begin{takeaway}
  \end{takeaway}
  \slidesources{Huyen2022}
\end{frame}

% TEMPLATE: captioned code panel. A model card is a document with fixed
% fields, so the honest slot is a panel showing those fields rather than a
% sentence about them. The caption is the artifact the heading names; the panel
% takes the book's field list and the question each field is there to answer.
\begin{frame}{Create model cards}
  \begin{tdcodet}{Model card}
  \end{tdcodet}
  \slidesources{Huyen2022}
\end{frame}

% TEMPLATE: withmargin. The heading asks for a process and not a list of peers,
% so the main column takes what that process has to cover and the failure mode
% the book names, fixes bolted on after the fact. Margin holds the tooling and
% the audits the book points at.
\begin{frame}{Establish processes for mitigating biases}
  \begin{withmargin}
    \begin{tdmain}
    \end{tdmain}
    \begin{tdmargin}
    \end{tdmargin}
  \end{withmargin}
  \slidesources{Huyen2022}
\end{frame}

% TEMPLATE: three block row, levelled by the equal height group, and a break
% from the column layout before it. The book answers this heading with named
% places to follow the field, so one box per kind of source rather than a run
% of names inside a sentence. Boxes stay untitled until those names are in.
\begin{frame}{Stay up-to-date on responsible AI}
  \begin{columns}[T,onlytextwidth]
    \begin{column}{0.31\textwidth}
      \begin{tdblock}[equal height group=rai]
      \end{tdblock}
    \end{column}
    \begin{column}{0.31\textwidth}
      \begin{tdblock}[equal height group=rai]
      \end{tdblock}
    \end{column}
    \begin{column}{0.31\textwidth}
      \begin{tdblock}[equal height group=rai]
      \end{tdblock}
    \end{column}
  \end{columns}
  \slidesources{Huyen2022}
\end{frame}

% ----------------------------------------------------------------------
\subsection{Summary}

% TEMPLATE: statement, kept bare. A summary is a closing line and not a slide
% of bullets, so it takes the loud centred form and nothing else. The main line
% is the book's own summing up of the human side; the substatement is the
% consequence it carries out of the chapter.
\begin{frame}{Summary}
  \begin{statement}
    \substatement{}
  \end{statement}
  \slidesources{Huyen2022}
\end{frame}

% ----------------------------------------------------------------------
% This chapter's own references. The refsection opened above scopes the
% citation numbering to this chapter, so chapter_pdfs/<this>.pdf resolves
% its own [n] markers instead of pointing at a list it does not contain.
\begin{frame}{References}
  \printbibliography[heading=none]
\end{frame}
\end{refsection}
